% 摘要与关键词
\begin{center}
\begin{minipage}{0.92\textwidth}
\small
\noindent\textbf{摘要}\quad 一个信号“能预测”与“能赚钱”之间的距离，通常被归到成本账上。这笔账只算对了一半。成本是拦路者之一，但完整的落差里还有另外两段：观测它的时间尺度，以及给它定价的参照系。同一批数据、同一个信号，日频调仓亏钱而数日调仓赚钱；在一个基准下显著的 $\alpha$，换一个基准就消失。这两类现象在实务中不新鲜，却很少被写进同一套可证明、可证伪的语言。本文尝试给出这套语言。

全文的核心是把“可交易缝隙”定义为一个相对于观测装置的存在性问题。观测装置由时间尺度与参照系两个维度刻画，缝隙的存在性则是“市场过程 $\times$ 观测装置”的联合属性。这一点有两组结果作支撑。其一（时间尺度）：在信息按指数衰减、单次交易付出比例成本的最小模型里，可提取的盈利空间被一对相反的约束夹成一个尺度窗口——网格太密，检测被观测噪声挡住；网格太疏，信息已经衰减。窗口的位置与宽度由信号的衰减速率决定。其二（参照系）：任何 $\alpha$ 都相对于一组比较工具（基准）而定义；参照系扩张时，可观测的缝隙恰好少掉“新增工具中与旧参照系正交、且有非零溢价”的那一部分，在完备参照系下归零。把这两组结果与成本、选择偏差校正合起来，得到一条四道核验的判据；实务中常用的“三道门”与统计门槛，可以解释为它的操作版本。

实证部分用本地通达信日线与沪深北 400 只个股的八年样本（2019--2026），对 8 个价量信号做尺度扫描：多空口径下 4 个候选出现清晰的内点峰值，位置集中在 2--5 个交易日；最慢的信号（非流动性）窗口右端全开，成本后夏普从 1 日尺度上的 $1.01$ 升到 20 日尺度上的 $1.13$；最快的信号（1 日反转）在全部尺度上被成本压制。这些曲线为尺度窗口提供了直接证据，也给出“什么频率上存在什么缝隙”的定量轮廓。文末如实列出判据的失效边界：时点安全的纪律、容量与拥挤对衰减速率的内生反馈，以及方向选择中无法完全消除的样本内成分。

\vspace{0.5em}
\noindent\textbf{关键词}\quad 市场有效性；交易成本；信息衰减；观测尺度；参照系；多重检验
\end{minipage}
\end{center}

\newpage